\documentclass[11pt]{article}
\usepackage[margin=1in]{geometry}
\usepackage{amsmath,amssymb,amsthm}
\usepackage{booktabs}
\usepackage{url}

\newtheorem{theorem}{Theorem}
\newtheorem{corollary}{Corollary}
\newtheorem{conjecture}{Conjecture}
\theoremstyle{definition}
\newtheorem{definition}{Definition}

\title{Disproof of TLMC Conjecture 00000000124:\\
the least missing value in the orbit of $2$ equals $q$, not $O((\log q)^2)$}
\author{TLMC falsification package}
\date{September 29, 2026}

\begin{document}
\maketitle

\begin{abstract}
Conjecture 00000000124 asserts that if $q$ is an odd prime for which $2$ is a
primitive root, then the least positive integer $m(q)$ missing from the orbit
$\{2^n \bmod q : n \ge 0\}$ satisfies $m(q) = O((\log q)^2)$.  This is false:
the hypothesis forces the orbit to be all of $\{1,\dots,q-1\}$, so
$m(q) = q$ identically, and $q > (\log q)^2$ for every $q \ge 3$.  The claim
therefore fails with constant $C=1$ at \emph{every} admissible prime (e.g.\
$m(3)=3>1.2069=(\ln 3)^2$ and $m(101)=101 \gg 21.2993=(\ln 101)^2$), and for
every fixed $C>0$ at every admissible prime $q>e^{6C}$.  All instances are
confirmed by independent recomputation and by machine-checked Lean~4
certificates using no axioms and no \texttt{sorry}.
\end{abstract}

\section{The conjecture and the attack}

\begin{conjecture}[TLMC 00000000124]
Let $q$ be an odd prime such that $2$ is a primitive root modulo $q$ (i.e.\
the multiplicative order satisfies $\operatorname{ord}_q(2)=q-1$).  Let $m(q)$
be the least positive integer that does not occur in the orbit
$\{2^n \bmod q : n \ge 0\}$, residues taken in $\{0,1,\dots,q-1\}$.  Then
$m(q) = O((\log q)^2)$.
\end{conjecture}

The attack is structural and requires no computation: the hypothesis is so
strong that it determines $m(q)$ exactly.

\section{The exact value of $m(q)$}

\begin{theorem}\label{thm:main}
Let $q$ be an odd prime with $\operatorname{ord}_q(2)=q-1$.  Then the orbit of
$2$ modulo $q$, written with least positive residues, is exactly
$\{1,2,\dots,q-1\}$, and the least missing positive integer is
\[
m(q) = q .
\]
\end{theorem}

\begin{proof}
Since $q$ is odd, $\gcd(2,q)=1$, so no power of $2$ is divisible by $q$: the
residues $r_n := 2^n \bmod q$ all lie in $\{1,\dots,q-1\}$.

The values $r_0, r_1, \dots, r_{q-2}$ are pairwise distinct.  Indeed, if
$2^i \equiv 2^j \pmod q$ with $0 \le i < j \le q-2$, then multiplying by the
inverse of $2^i$ gives $2^{\,j-i} \equiv 1 \pmod q$ with $0 < j-i < q-1$,
contradicting $\operatorname{ord}_q(2) = q-1$.

So the orbit contains $q-1$ distinct elements of the $(q-1)$-element set
$\{1,\dots,q-1\}$, hence equals it: $\{r_0,\dots,r_{q-2}\} = \{1,\dots,q-1\}$
(and $r_n$ for $n \ge q-1$ merely repeats it, since $2^{q-1}\equiv 1$).

Consequently every positive integer $k \le q-1$ occurs in the orbit (as the
residue $k$ itself), while $q$ does not: $q \equiv 0 \pmod q$ and $0$ is not a
power of $2$ modulo $q$.  The least missing positive integer is therefore
$q$.  \qed
\end{proof}

\section{The gap against the logarithmic bound}

\begin{theorem}\label{thm:gap}
For every real $x \ge 1$, $(\ln x)^2 < x$.
\end{theorem}

\begin{proof}
Put $t=\ln x \ge 0$.  We show $e^{t/2} > t$, i.e.\ $\sqrt{x} > \ln x$.
The tangent bound $e^u \ge 1+u$ (all Taylor coefficients are nonnegative) gives,
for $0 \le t < 2$: $e^{t/2} \ge 1 + t/2 > t$.  For $t \ge 2$, applying the same
bound at $u=(t-2)/2 \ge 0$ gives
$e^{t/2} = e \cdot e^{(t-2)/2} \ge e\,(1+\tfrac{t-2}{2}) = \tfrac{e}{2}\,t > t$,
because $e > 2$.  \qed
\end{proof}

\begin{corollary}\label{cor:refute}
For every odd prime $q$ with $\operatorname{ord}_q(2)=q-1$,
\[
m(q) = q > (\log q)^2
\]
with the constant $C=1$; already $m(3)=3 > 1.2069\ldots=(\ln 3)^2$ and
$m(101)=101 \gg 21.2993\ldots=(\ln 101)^2$.  More generally, for every fixed
$C>0$: since $e^{L} \ge L^3/6$ for $L\ge0$ (Taylor series with nonnegative
terms), every $q$ with $L=\ln q > 6C$ satisfies
$q = e^{L} \ge \tfrac{L^3}{6} > C L^2 = C(\log q)^2$, hence
$m(q) > C(\log q)^2$.
\end{corollary}

\section{Scope and boundary of the refutation}
\label{sec:boundary}

The following statements are \textbf{unconditional}:
\begin{itemize}
\item $m(q)=q$ exactly, for every admissible $q$ (Theorem~\ref{thm:main});
\item $m(q) > (\log q)^2$, i.e.\ the conjecture fails with constant $C=1$, for
      every admissible $q$ without exception (Corollary~\ref{cor:refute});
\item for each fixed $C \le $ any prescribed value, every admissible
      $q > e^{6C}$ violates the proposed bound.
\end{itemize}

The full refutation of the $O(\cdot)$ statement for \emph{arbitrary} constants
$C$ additionally requires that admissible primes be unbounded.  By
Theorem~\ref{thm:main}, $m(q)/(\log q)^2 = q/(\log q)^2 \to \infty$ along any
unbounded sequence of admissible primes; the infinitude of primes with $2$ as
primitive root is Artin's primitive-root conjecture for $a=2$, which is open in
general but provable under the Generalized Riemann Hypothesis
(Hooley, 1967) and supported by the expected positive density
($\approx 0.374$; OEIS A001122).  Note that if the admissible set were finite,
the $O$-statement would hold vacuously for large $q$; thus the exact identity
$m(q)=q$ is the unconditional core of the disproof, and it already destroys
the conjecture's polylogarithmic content: the true growth is \emph{linear} in
$q$, not $(\log q)^2$.  This boundary is stated explicitly in
\texttt{README.md}.

\section{Computational verification}

An independent recomputation (\texttt{reproduce.py}, pure standard library)
confirms the verdict's spot checks and the general picture for all $22$
admissible primes below $200$: the orbit is $\{1,\dots,q-1\}$, $m(q)=q$, and
$m(q) > (\ln q)^2$ throughout, with the ratio $m(q)/(\ln q)^2$ growing
steadily.

\begin{center}
\begin{tabular}{rrrr@{\hspace{2em}}rrrr}
\toprule
$q$ & $m(q)$ & $(\ln q)^2$ & ratio & $q$ & $m(q)$ & $(\ln q)^2$ & ratio\\
\midrule
3   & 3   & 1.207  & 2.49 & 83  & 83  & 19.526 & 4.25\\
5   & 5   & 2.590  & 1.93 & 101 & 101 & 21.299 & 4.74\\
11  & 11  & 5.750  & 1.91 & 107 & 107 & 21.835 & 4.90\\
13  & 13  & 6.579  & 1.98 & 131 & 131 & 23.768 & 5.51\\
19  & 19  & 8.670  & 2.19 & 139 & 139 & 24.349 & 5.71\\
29  & 29  & 11.339 & 2.56 & 149 & 149 & 25.039 & 5.95\\
37  & 37  & 13.039 & 2.84 & 163 & 163 & 25.946 & 6.28\\
53  & 53  & 15.763 & 3.36 & 173 & 173 & 26.556 & 6.51\\
59  & 59  & 16.626 & 3.55 & 179 & 179 & 26.909 & 6.65\\
61  & 61  & 16.899 & 3.61 & 181 & 181 & 27.024 & 6.70\\
67  & 67  & 17.679 & 3.79 & 197 & 197 & 27.912 & 7.06\\
\bottomrule
\end{tabular}
\end{center}

The two verdict spot checks are reproduced exactly: for $q=3$,
$m=3 > (\ln 3)^2 \approx 1.2069$; for $q=101$, $m=101 \gg (\ln 101)^2 \approx
21.2993$.

\section{Machine-checked Lean 4 certificates}

The directory \texttt{lean4/} contains a self-contained Lean~4 project
(toolchain \texttt{leanprover/lean4:v4.33.1}, no dependencies, core Lean only).
For each admissible prime $q<200$ it proves, by \texttt{decide} over $\mathbb{N}$
arithmetic or by the two elementary lemmas of the file:
\begin{itemize}
\item \texttt{pr\_q}: $2$ is a primitive root mod $q$ --- Fermat's congruence
      $2^{q-1}\equiv 1$ together with $2^{(q-1)/f}\not\equiv 1$ for each prime
      $f \mid q-1$;
\item \texttt{cover\_q}: every $1 \le k < q$ occurs in the orbit
      $\{2^n \bmod q\}$, via explicit discrete-log witness lists checked by
      \texttt{dlogOk} and the semantic lemma \texttt{dlogOk\_implies};
\item \texttt{lm\_q}: \texttt{IsLeastMissing q q}, i.e.\ the least positive
      integer missing from the orbit is exactly $q$ (so $m(q)=q$);
\item \texttt{pw\_q}: $2^{\ell} \le q < 2^{\ell+1}$ for
      $\ell=\lfloor\log_2 q\rfloor$ and $\ell^2 < q$ --- the concrete gap
      $m(q)=q > (\log_2 q)^2$.
\end{itemize}
The flagship \texttt{attack\_101} bundles the verdict's $q=101$ check:
primitive-root certificate, $m(101)=101$, and
$6^2=36=(\lfloor\log_2 101\rfloor)^2 < 101$.

\texttt{Check.lean} runs \texttt{\#print axioms} on all $138$ theorems; every
single one reports \emph{does not depend on any axioms} --- zero
\texttt{sorry}, zero \texttt{propext}/\texttt{Quot.sound}/%
\texttt{Classical.choice}, i.e.\ a fully constructive, axiom-free audit.

\section{Reproducing}

\begin{verbatim}
python3 reproduce.py                 # independent recomputation (stdlib only)
cd lean4 && lake build               # builds Main.lean (Lean 4.33.1)
lake env lean Check.lean             # axiom audit: 138 x "does not depend on any axioms"
tectonic main.tex --outdir build     # this document
\end{verbatim}

\begin{thebibliography}{9}
\bibitem{artin} E.~Artin, \emph{\"Uber eine neue Art von L-Reihen},
Abh.\ Math.\ Sem.\ Univ.\ Hamburg 3 (1927), 89--108.
\bibitem{hooley} C.~Hooley, \emph{On Artin's conjecture}, J.\ Reine Angew.\
Math.\ 225 (1967), 209--220.
\bibitem{oeis} OEIS Foundation, \emph{A001122: Primes with primitive root 2},
\url{https://oeis.org/A001122}.
\end{thebibliography}

\end{document}
